\section{Concurrencia y transferencias}

El nucleo transaccional del mini banco vive en memoria, sin SGBD, como exige el
enunciado. La integridad del dinero descansa en dos decisiones de diseno: el
saldo se representa siempre en centavos como \codigo{long} (nunca \codigo{double})
y cada transferencia toma los candados de las dos cuentas en un orden total para
evitar interbloqueos. Los componentes implicados son
\codigo{model/Cuenta.java} (el dato y su candado) y
\codigo{data/CuentaRepository.java} (la logica concurrente y el invariante).

\subsection{Dinero en centavos: exactitud y conservacion}

La clase \codigo{Cuenta} guarda el saldo en un \codigo{long} de centavos. El uso
de aritmetica entera elimina los errores de redondeo del punto flotante, de modo
que sumas y restas son exactas y el invariante de conservacion del total se
mantiene al centavo. La salida JSON formatea el saldo directo desde el
\codigo{long} con \codigo{appendSaldo(StringBuilder)} (sin crear un
\codigo{BigDecimal} por GET); \codigo{saldoDecimal()} se conserva como conversion
exacta de respaldo, escalando el entero por $10^{-2}$.

\begin{lstlisting}[language=Java,caption={Saldo entero y conversion exacta en Cuenta.java}]
private volatile long saldoCentavos;

public long getSaldoCentavos() { return saldoCentavos; }
public void setSaldoCentavos(long c) { this.saldoCentavos = c; }

/** Saldo como decimal exacto (centavos / 100) para el JSON de salida. */
public BigDecimal saldoDecimal() {
    return BigDecimal.valueOf(saldoCentavos, 2);
}
\end{lstlisting}

Cada \codigo{Cuenta} expone ademas un \codigo{ReentrantLock} propio, que es la
unidad de exclusion mutua usada por las transferencias. Tener un candado por
cuenta (en vez de uno global) permite que transferencias sobre cuentas disjuntas
avancen en paralelo.

\begin{lstlisting}[language=Java,caption={Candado por cuenta en Cuenta.java}]
// Lock por cuenta para transferencias concurrentes (se toma en orden de id).
public final ReentrantLock lock = new ReentrantLock();
\end{lstlisting}

\subsection{Estructura concurrente del repositorio}

\codigo{CuentaRepository} es un singleton (\codigo{CuentaRepository.get()}) que
almacena las cuentas en un \codigo{ConcurrentHashMap} indexado por id. El
\codigo{AtomicLong secuencia} es el numero de secuencia monotono que ordena las
transacciones (base de la replicacion y del log en GCS). El conteo de commits
\codigo{totalTransferencias} es un \codigo{LongAdder}: en el hot path hace
\codigo{increment()} sin CAS global (\codigo{sum()} solo lo lee \codigo{/stats}),
quitando uno de los dos puntos de serializacion global por transferencia. El
campo \codigo{ultimaTxId} es \codigo{volatile} para publicar el ultimo seq sin
candado adicional. Por ultimo, el \codigo{StampedLock snapshotLock} coordina los
movimientos con el escaneo del total: cada transferencia lo toma en modo
\emph{compartido} (corren en paralelo entre si) y la suma del total lo toma en
modo \emph{exclusivo}, para que el agregado nunca se lea a mitad de un movimiento.

\begin{lstlisting}[language=Java,caption={Estado concurrente de CuentaRepository.java}]
private final ConcurrentHashMap<Integer, Cuenta> cuentas = new ConcurrentHashMap<>();

private final AtomicLong secuencia = new AtomicLong(0);
private final LongAdder totalTransferencias = new LongAdder();
private volatile long ultimaTxId = 0;
private final RegistroTransacciones registro = new RegistroTransacciones();

// Serializa secuencia+registro.agregar+observador: lock mas interno (sin deadlock).
private final ReentrantLock seqLock = new ReentrantLock();

// Snapshot: los movimientos toman readLock (compartido); el escaneo del total,
// writeLock (exclusivo). Asi saldoTotalCentavos() es una foto consistente.
private final StampedLock snapshotLock = new StampedLock();
\end{lstlisting}

El invariante de correctitud es que la suma de todos los saldos permanece
constante ante cualquier secuencia de transferencias. El metodo
\codigo{saldoTotalCentavos()} lo materializa recorriendo las cuentas y sumando
sus centavos. Para que ese recorrido sea una \emph{foto consistente} y nunca
observe un movimiento a medias (una cuenta ya debitada y su contraparte aun sin
acreditar), toma el \codigo{writeLock} exclusivo del \codigo{snapshotLock}: asi
ningun \codigo{transferir()} (que sostiene el \codigo{readLock} compartido) puede
estar en su tramo debito+credito durante la suma. Sigue siendo una suma real
sobre los saldos --no un contador incremental--, de modo que tambien detecta un
descuadre (corrupcion o replica mal aplicada); es la propiedad que verifica la
prueba \codigo{detectaDescuadre}. El costo es minimo: el escaneo es raro (cada
\textasciitilde2\,s en \codigo{/stats}) y solo entonces pausa los movimientos
unos milisegundos.

\begin{lstlisting}[language=Java,caption={Foto consistente del invariante en CuentaRepository.java}]
/** Suma consistente de todos los saldos (writeLock: ningun movimiento a medias). */
public long saldoTotalCentavos() {
    long stamp = snapshotLock.writeLock();
    try {
        long total = 0;
        for (Cuenta c : cuentas.values()) total += c.getSaldoCentavos();
        return total;
    } finally {
        snapshotLock.unlockWrite(stamp);
    }
}
\end{lstlisting}

\subsection{transferir(): atomicidad sin interbloqueo}

\codigo{CuentaRepository.transferir()} es la operacion central. Primero rechaza
los casos triviales (misma cuenta, monto no positivo) y la ausencia de cualquiera
de las cuentas. Acto seguido aplica la clave anti-deadlock: ordena los dos
candados por id (el menor primero) de manera que dos hilos que toquen el mismo par
de cuentas siempre los adquieran en el mismo orden, eliminando la espera
circular. Bajo ambos candados valida el saldo y mueve el monto; luego, dentro de
un tercer candado \codigo{seqLock} (el mas interno), asigna el seq, registra la
\codigo{Transaccion} y avisa al observador en un mismo tramo serializado, de modo
que el orden de insercion en el registro coincida con el orden de seq (evita
perder transacciones en la replicacion en vivo). El bloque \codigo{finally}
garantiza la liberacion en orden inverso.

\begin{lstlisting}[language=Java,caption={Bloqueo ordenado y commit en CuentaRepository.java}]
// Bloqueo en orden de id (menor primero) para no provocar deadlock. El readLock
// del snapshot envuelve el par debito+credito (se toma ANTES de los locks de
// cuenta y se suelta DESPUES) para que el escaneo del total no caiga en medio.
Cuenta primero = origenId < destinoId ? origen : destino;
Cuenta segundo = origenId < destinoId ? destino : origen;
long stamp = snapshotLock.readLock();
primero.lock.lock();
segundo.lock.lock();
try {
    if (origen.getSaldoCentavos() < montoCentavos) {
        return Resultado.de(Estado.SALDO_INSUFICIENTE);
    }
    origen.setSaldoCentavos(origen.getSaldoCentavos() - montoCentavos);
    destino.setSaldoCentavos(destino.getSaldoCentavos() + montoCentavos);
    // Tramo serializado: orden de insercion en registro == orden de seq.
    seqLock.lock();
    try {
        long seq = secuencia.incrementAndGet();
        totalTransferencias.increment();
        ultimaTxId = seq;
        Transaccion tx = new Transaccion(seq, origenId, destinoId, montoCentavos);
        registro.agregar(tx);
        observador.accept(tx);
        return Resultado.ok(seq);
    } finally {
        seqLock.unlock();
    }
} finally {
    segundo.lock.unlock();
    primero.lock.unlock();
    snapshotLock.unlockRead(stamp);
}
\end{lstlisting}

\begin{nota}
Como cada cuenta cede el mismo monto que recibe la otra, el delta neto sobre el
total es cero: por eso \codigo{saldoTotalCentavos()} no cambia tras un
\codigo{transferir()} exitoso.
\end{nota}

El observador y el registro son los puntos de enganche con replicacion y
durabilidad. \codigo{setObservador()} instala un \codigo{Consumer<Transaccion>}
(lo pone el lider) al que se notifica cada commit para propagar la transaccion a
las replicas y al log durable; \codigo{setConteoStorage()} inyecta el numero de
transacciones presentes en Cloud Storage, consultable con \codigo{txEnStorage()}.

\begin{lstlisting}[language=Java,caption={Puntos de enganche en CuentaRepository.java}]
public void setObservador(Consumer<Transaccion> o) { this.observador = o; }
public void setConteoStorage(LongSupplier s) { this.conteoStorage = s; }
public long txEnStorage() { return conteoStorage.getAsLong(); }
\end{lstlisting}

La replica aplica los cambios sin revalidar: \codigo{aplicarReplica()} confia en
que el lider ya valido el saldo, por lo que solo mueve el monto, registra la tx y
avanza la secuencia local al maximo entre la actual y \codigo{t.seq()} mediante
\codigo{accumulateAndGet(t.seq(), Math::max)}. Reusa el mismo bloqueo ordenado por
id para no romper la coherencia frente a transferencias locales concurrentes.

\begin{lstlisting}[language=Java,caption={Aplicacion idempotente en la replica, CuentaRepository.java}]
origen.setSaldoCentavos(origen.getSaldoCentavos() - t.montoCentavos());
destino.setSaldoCentavos(destino.getSaldoCentavos() + t.montoCentavos());
registro.agregar(t);
totalTransferencias.increment();
ultimaTxId = t.seq();
secuencia.accumulateAndGet(t.seq(), Math::max);
\end{lstlisting}

\subsection{Errores de negocio y mapeo a codigos REST}

\codigo{transferir()} no lanza excepciones de negocio: devuelve un
\codigo{Resultado} con un \codigo{Estado} de la enumeracion. La capa HTTP,
\codigo{server/TransferHandler.java}, traduce ese estado a un codigo de respuesta
y un cuerpo JSON. El exito devuelve el seq asignado; los fallos de negocio se
agrupan en \codigo{400} salvo la cuenta inexistente, que es \codigo{404}.

\begin{tablaglab}{L{4.5cm} C{2cm} L{6cm}}
\thead{Estado} & \thead{HTTP} & \thead{Significado} \\ \midrule
OK                & 200 & Commit; el cuerpo incluye el \codigo{seq}. \\
NO\_ENCONTRADA    & 404 & Origen o destino no existe. \\
MISMA\_CUENTA     & 400 & Origen y destino iguales. \\
MONTO\_INVALIDO   & 400 & Monto menor o igual a cero. \\
SALDO\_INSUFICIENTE & 400 & El origen no cubre el monto. \\
\end{tablaglab}

\begin{lstlisting}[language=Java,caption={Mapeo de Estado a respuesta REST en TransferHandler.java}]
CuentaRepository.Resultado r = repo.transferir(origen, destino, centavos);
switch (r.estado) {
    case OK:                return Respuesta.json(200, "{\"status\":\"ok\",\"seq\":" + r.seq + "}");
    case NO_ENCONTRADA:     return Respuesta.error(404, "Cuenta no encontrada");
    case MISMA_CUENTA:      return Respuesta.error(400, "Origen y destino iguales");
    case MONTO_INVALIDO:    return Respuesta.error(400, "Monto invalido");
    case SALDO_INSUFICIENTE: return Respuesta.error(400, "Saldo insuficiente");
    default:                return Respuesta.error(500, "Error interno");
}
\end{lstlisting}

\begin{advertencia}
La validacion de saldo ocurre dentro de los candados, no antes: comprobar el
saldo fuera de la seccion critica abriria una condicion de carrera (dos retiros
concurrentes podrian pasar la verificacion y dejar saldo negativo).
\end{advertencia}
